\subsection{Descent by an étale and surjective map}

Let X and Y be topological spaces, and let f be a continuous map from X to Y. We retain the notation of the preceding section.

THEOREM 2. --- Suppose that every point of Y has a neighborhood over which there exists a continuous section of the map f. Then the groupoid morphism \(\varpi'(f)\) from \(\operatorname{Coeg}(f)\) into \(\varpi(Y)\) is an isomorphism.

By hypothesis, there exists a covering \((\mathrm{U}_{j})_{j\in\mathrm{J}}\) of Y by open sets and, for every \(j\in J\), a continuous section \(s_{j}\) of \(f_{U_{j}}\) .

If \(c\) is a path in X, we will denote by \([c]\) its strict homotopy class in \(\varpi(X)\) and by \(\{c\}\) the image of \([c]\) in \(\operatorname{Coeg}(f)\) under the groupoid morphism \(\gamma\) . If \(c\) and \(c'\) are two paths in X, \(\{c\}\) and \(\{c'\}\) are composable in \(\operatorname{Coeg}(f)\) if and only if the paths \(f \circ c\) and \(f \circ c'\) in Y are juxtaposable.

Let \(c'\) be a path in Y. By Lemma 4 of III, p.~272, applied to the compact space I and the covering \(((c')^{-1}(\mathrm{U}_j))_{j\in \mathrm{J}}\) of I, there exists an integer \(n\) such that, for every integer \(k\) satisfying \(1\leqslant k\leqslant n\), the image of the interval \([\frac{k - 1}{n},\frac{k}{n}]\) under \(c'\) is contained in an open set \(\mathrm{U}_{j(k)}\) . For every integer \(k\) , \(1\leqslant k\leqslant n\) , let \(c_k'\) be the path in Y defined by \(s\mapsto c'(\frac{k + s - 1}{n})\) ; we have \([c'] = [c_1'][c_2']\ldots [c_n']\) (cf. III, p.~291, Remark 1), and \(c'\) is the path denoted \(c_1' * c_2' * \dots * c_n'\) . For all \(k\in \{1,\ldots ,n\}\) , denote by \(c_k\) the path \(s_{j(k)}\circ c_k'\) in X and set \(\{c_k\} = \gamma ([c_k])\) . Since for every \(k\) the paths \(c_{k - 1}' = f\circ c_{k - 1}\) and \(c_k' = f\circ c_k\) are juxtaposable, the sequence \((\{c_1\} ,\ldots ,\{c_n\})\) is composable in \(\operatorname{Coeg}(f)\). By construction,

\[
\begin{array}{r l} & {\varpi^ {\prime} (f) (\{c _ {1} \} \ldots \{c _ {n} \}) = \varpi^ {\prime} (f) (\{c _ {1} \}) \ldots \varpi^ {\prime} (f) (\{c _ {n} \})} \\ & {\qquad = \varpi (f) ([ c _ {1} ]) \ldots \varpi (f) ([ c _ {n} ])} \\ & {\qquad = [ f \circ c _ {1} ] \ldots [ f \circ c _ {n} ]} \\ & {\qquad = [ c _ {1} ^ {\prime} ] * \dots * [ c _ {n} ^ {\prime} ] = [ c ^ {\prime} ],} \end{array}
\]

which proves that the groupoid morphism \(\varpi'(f)\) is surjective.

Let \(u\) and \(v\) be arrows of \(\operatorname{Coeg}(f)\) . Since the groupoid \(\operatorname{Coeg}(f)\) is generated by the image of \(\varpi(\mathrm{X})\) (II, p.~200, corollary), there exist finite sequences \((c_1, \ldots, c_n)\) and \((d_1, \ldots, d_n)\) of paths in X such that one has \(u = \{c_1\} \ldots \{c_n\}\) and \(v = \{d_1\} \ldots \{d_n\}\) . The paths \((f \circ c_1, \ldots, f \circ c_n)\) in Y are then juxtaposable and one has \(\varpi'(f)(u) = [(f \circ c_1)] \ldots [(f \circ c_n)]\) ; likewise, \(\varpi'(f)(v) = [(f \circ d_1)] \ldots [(f \circ d_n)]\) .

Suppose that we have \(\varpi'(f)(u) = \varpi'(f)(v)\) . There then exists a strict homotopy \(\sigma\) connecting \((f \circ c_1) * \cdots * (f \circ c_n)\) to \((f \circ d_1) * \cdots * (f \circ d_n)\) . By Lemma 4 of III, p.~272, applied to the compact space \(\mathbf{I} \times \mathbf{I}\) and the covering \((\bar{\sigma}^{-1}(\mathrm{U}_i))_{i \in \mathrm{I}}\) of \(\mathbf{I} \times \mathbf{I}\) , there exists an integer \(m \geqslant 1\) such that, for every pair of integers \((j,k)\) satisfying \(1 \leqslant j \leqslant m\) and \(1 \leqslant k \leqslant m\), the image of \([\frac{j-1}{m}, \frac{j}{m}] \times [\frac{k-1}{m}, \frac{k}{m}]\) under \(\sigma\) is contained in an open set \(U_{i(j,k)}\) of the covering \((\mathrm{U}_{i})_{i \in \mathrm{I}}\) .

Every path c in X is of the form \(c_{1}*\cdots*c_{m}\) , where \(c_{k}\) is the path \(t\mapsto c(\frac{k-1+t}{m})\) . If necessary, replacing the integers m and n by their product mn, one may therefore suppose that m=n.

For every pair \((j,k)\) of integers of \(\{1,\ldots ,n\}\) and every pair \((s,t)\in\) \(\mathbf{I}\times \mathbf{I}\) , set

\[
\sigma_ {j, k} (s, t) = s _ {i (j, k)} \circ \sigma \left(\frac {s + j - 1}{n}, \frac {t + k - 1}{n}\right).
\]

For \(t \in \mathbf{I}\) , also set \(h_{j,k}^{0}(t) = \sigma_{j,k}(t,0)\) , \(h_{j,k}^{1}(t) = \sigma_{j,k}(t,1)\) , \(v_{j,k}^{0}(t) = \sigma_{j,k}(0,t)\) and \(v_{j,k}^{1}(t) = \sigma_{j,k}(1,t)\) . By Lemma 1 of III, p.~295, the paths \(h_{j,k}^{0} * v_{j,k}^{1}\) and \(v_{j,k}^{0} * h_{j,k}^{1}\) are strictly homotopic, whence the relation

\[
[ h _ {j, k} ^ {0} ] [ v _ {j, k} ^ {1} ] = [ v _ {j, k} ^ {0} ] [ h _ {j, k} ^ {1} ]\tag{2}
\]

on \(\varpi(\mathrm{X})\) , for every pair \((j,k) \in \{1, \ldots, n\}^2\) . On the other hand, for every pair of integers \((j,k)\) , with \(2 \leqslant j \leqslant n\) and \(1 \leqslant k \leqslant n\) , we have \(f \circ v_{j,k}^0 = f \circ v_{j-1,k}^1\) , whence the relation

\[
\{v _ {j, k} ^ {0} \} = \{v _ {j - 1, k} ^ {1} \}\tag{3}
\]

in Coeg(f). Likewise, for every pair of integers \((j,k)\) such that \(1 \leqslant j \leqslant n\) and \(2 \leqslant k \leqslant n\) , we have

\[
\{h _ {j, k} ^ {0} \} = \{h _ {j, k - 1} ^ {1} \}.\tag{4}
\]

For \(j \in \{1, \ldots, n\}\) the paths \(f \circ c_j\) and \(f \circ h_{j,1}^0\) coincide. By definition of the coequalizer Coeg(f), one therefore has \(\{c_j\} = \{h_{j,1}^0\}\) . Likewise, for every \(j \in \{1, \ldots, n\}\) , \(\{d_j\} = \{h_{j,n}^1\}\) . Consequently, one has the relations

\[
u = \{h _ {1, 1} ^ {0} \} \dots \{h _ {n, 1} ^ {0} \}, \quad v = \{h _ {1, n} ^ {1} \} \dots \{h _ {n, n} ^ {1} \}.
\]

By Lemma 3 below, one has

\[
\{h _ {1, 1} ^ {0} \} \dots \{h _ {1, n} ^ {0} \} \{v _ {1, n} ^ {1} \} \dots \{v _ {n, n} ^ {1} \} = \{v _ {1, 1} ^ {0} \} \dots \{v _ {n, 1} ^ {0} \} \{h _ {n, 1} ^ {1} \} \dots \{h _ {n, n} ^ {1} \}.
\]

Since the paths \(t \mapsto \sigma(0,t)\) and \(t \mapsto \sigma(1,t)\) are constant, one has, for \(1 \leqslant k \leqslant n\) , \(\{v_{1,k}^{0}\} = e_{a}\) and \(\{v_{n,k}^{1}\} = e_{b}\) where \(a,b \in Y\) are the source and the target of the arrows \(u\) and \(v\) of Coeg( \(f\) ). It follows that one has the equality

\[
\{h _ {1, 1} ^ {0} \} \dots \{h _ {1, n} ^ {0} \} = \{h _ {n, 1} ^ {0} \} \dots \{h _ {n, n} ^ {1} \},
\]

that is, u = v.

Lemma 3. --- Let G be a groupoid and let p, q be integers \(\geqslant 1\) . For every pair \((j, k)\) of integers such that \(0 \leqslant j \leqslant p\) and \(0 \leqslant k \leqslant q\) , let \(x_{j,k}\) be a vertex of G; for every pair \((j, k)\) such that \(1 \leqslant j \leqslant p\) and \(0 \leqslant k \leqslant q\) , let \(h_{j,k}\) be an arrow of G connecting \(x_{j-1,k}\) to \(x_{j,k}\) ; for every pair \((j, k)\) such that \(0 \leqslant j \leqslant p\) and \(1 \leqslant k \leqslant q\) , let \(v_{j,k}\) be an arrow of G connecting \(x_{j,k-1}\) to \(x_{j,k}\) .

Suppose that, for every pair \((j,k)\) of integers such that \(1 \leqslant j \leqslant p\) and \(1 \leqslant k \leqslant q\) , the arrows \(v_{j-1,k}\) and \(h_{j,k}\) are composable, as are the arrows \(h_{j,k-1}\) and \(v_{j,k}\) , and that one has \(v_{j-1,k}h_{j,k} = h_{j,k-1}v_{j,k}\) . Then,

\[
h _ {1, 0} h _ {2, 0} \dots h _ {p, 0} v _ {p, 1} v _ {p, 2} \dots v _ {p, q} = v _ {0, 1} v _ {0, 2} \dots v _ {0, q} h _ {1, q} h _ {2, q} \dots h _ {p, q}.
\]

First treat the particular case where q = 1 and prove the result by induction on p. If p = 1, the assertion to be proved is verified by hypothesis; suppose it is verified for p - 1; one then has

\[
h _ {1, 0} h _ {2, 0} \dots h _ {p, 0} v _ {p, 1} = h _ {1, 0} h _ {2, 0} \dots h _ {p - 1, 0} v _ {p - 1, 1} h _ {p, 1} = v _ {0, 1} h _ {1, 1} \dots h _ {p, 1}
\]

by the induction hypothesis, whence the relation for p.

Now prove the result by induction on q. It is true for q = 1 by what precedes; if it is true for q - 1, one then has

\[
h _ {1, 0} h _ {2, 0} \dots h _ {p, 0} v _ {p, 1} v _ {p, 2} \dots v _ {p, q} = v _ {0, 1} v _ {0, 2} \dots v _ {0, q - 1} h _ {1, q - 1} \dots h _ {p, q - 1} v _ {p, q}.
\]

By the case \(q = 1\) , we have

\[
h _ {1, q - 1} \dots h _ {p, q - 1} v _ {p, q} = v _ {0, q} h _ {1, q} \dots h _ {p, q},
\]

whence the desired relation.

Examples. --- 1) The theorem applies when the map f is étale and surjective.

\begin{enumerate}
\def\labelenumi{\arabic{enumi})}
\setcounter{enumi}{1}
\tightlist
\item
  It also applies when the space X is the sum space of a family \((V_{i})_{i\in\mathrm{I}}\) of subsets of Y whose interiors cover Y, and f is the map deduced from the canonical injections of each of the \(V_{i}\) in Y.
\end{enumerate}

